\documentclass[a4paper]{article}

\usepackage{graphicx}
\usepackage{amsmath,amssymb}
\usepackage{minted}
\usepackage{multirow}
\usepackage{float}
\usepackage{todonotes}
\usepackage{forest}
\usepackage{xstring}
\usepackage{xcolor}
\usepackage[most]{tcolorbox}
\usepackage{xparse}
\usepackage{natbib}
\usepackage{hyperref}

\usetikzlibrary{graphs,graphdrawing}
\usegdlibrary{layered}

% for external graphviz
\input{/home/donkarlo/Dropbox/repo/nd_ai_project/src/nd_ai/knoweldge_representation/language/formal/computer/programming/tex/latex/code_clips/external_graphviz.tex}

% for tags
\input{/home/donkarlo/Dropbox/repo/nd_ai_project/src/nd_ai/knoweldge_representation/language/formal/computer/programming/tex/latex/parts/control_sequence/kind/semtag/semtag.tex}

% for links
\input{/home/donkarlo/Dropbox/repo/nd_ai_project/src/nd_ai/knoweldge_representation/language/formal/computer/programming/tex/latex/code_clips/seehere.tex}

\title{Software, Platforms, and Frameworks for the AI Engineer Role}
\date{}
\author{Mohammad Rahmani}

\begin{document}
\maketitle

\section*{Target Role}
Production AI Engineer / AI Researcher focused on machine learning, GenAI, LLM systems, RAG, agentic workflows, Azure, Databricks, APIs, pipelines, deployment, monitoring, and business integration.

\section*{How to Use This Document}
This map translates the job description into concrete tools that can be mentioned in interviews, studied selectively, or used to plan a small portfolio project. The strongest match for this specific role is the combination of Python, Azure, Databricks, RAG and agent frameworks, and MLOps/monitoring.

\section{Production-ready AI Applications}

\begin{tcolorbox}[title={FastAPI}]
Use it to expose ML or LLM functionality as clean REST APIs in Python.
It is especially useful when a prototype must become a reliable internal service with typed inputs, validation, and automatic API documentation.
\end{tcolorbox}

\begin{tcolorbox}[title={Docker}]
Package the application, model code, dependencies, and runtime into reproducible containers.
This makes deployment to Azure, Kubernetes, or Databricks jobs more predictable and easier to maintain.
\end{tcolorbox}

\section{Traditional Machine Learning and Data Science}

\begin{tcolorbox}[title={scikit-learn}]
A strong default for classical ML, feature preprocessing, model selection, and baseline experiments.
It is ideal for tabular classification, regression, clustering, and quick comparison before moving to heavier deep learning solutions.
\end{tcolorbox}

\begin{tcolorbox}[title={PyTorch}]
A flexible deep learning framework for neural networks, custom training loops, and research-to-production workflows.
It is useful when the task needs deep models, embeddings, sequence models, or integration with modern LLM tooling.
\end{tcolorbox}

\section{GenAI and LLM-based Systems}

\begin{tcolorbox}[title={Azure OpenAI Service}]
Use it to access enterprise-managed LLMs through Azure security, networking, and governance.
It fits internal assistant systems where authentication, data protection, monitoring, and integration with Microsoft services matter.
\end{tcolorbox}

\begin{tcolorbox}[title={Hugging Face Transformers}]
Use it for open-source transformer models, embeddings, fine-tuning, and local or private experimentation.
It is helpful when model control, domain adaptation, or alternatives to closed cloud APIs are needed.
\end{tcolorbox}

\section{RAG Systems}

\begin{tcolorbox}[title={Azure AI Search}]
Use it as the retrieval layer for enterprise RAG with keyword, vector, semantic, hybrid, and multimodal search patterns.
It is a strong match when documents must be indexed securely and queried at scale inside Azure.
\end{tcolorbox}

\begin{tcolorbox}[title={LlamaIndex}]
Use it to build document ingestion, indexing, retrieval, and query pipelines for RAG applications.
It is especially good for connecting many data sources and turning unstructured internal documents into searchable LLM context.
\end{tcolorbox}

\section{Agentic AI and Workflow Automation}

\begin{tcolorbox}[title={LangGraph}]
Use it to build stateful, multi-step agents with controllable flows, tool calls, memory, retries, and human review points.
It is better than a simple chatbot when the AI must plan, call tools, validate outputs, and continue across steps.
\end{tcolorbox}

\begin{tcolorbox}[title={Microsoft Copilot Studio}]
Use it to create enterprise agents that connect to Microsoft 365, business data, and workflow actions.
It is useful when business units need practical automation without every use case becoming a full custom software project.
\end{tcolorbox}

\section{AI Services, APIs, and Integration into Existing Systems}

\begin{tcolorbox}[title={FastAPI}]
Use it as the interface layer between models, databases, frontends, and enterprise systems.
It lets you separate AI logic from application logic and expose stable endpoints for other teams.
\end{tcolorbox}

\begin{tcolorbox}[title={Azure Functions}]
Use it for serverless AI-related tasks such as document processing, scheduled inference, event-based automation, or lightweight workflow triggers.
It is useful when a full always-on service is unnecessary.
\end{tcolorbox}

\section{Modern AI Architecture and Model Orchestration}

\begin{tcolorbox}[title={LangChain}]
Use it to connect prompts, tools, retrievers, memory, models, and output parsers into LLM applications.
It is practical for quickly composing RAG and agent prototypes before hardening the most stable parts.
\end{tcolorbox}

\begin{tcolorbox}[title={Semantic Kernel}]
Use it when building LLM applications in the Microsoft ecosystem with planners, plugins, and structured orchestration.
It fits Azure-heavy environments where AI functions need to be connected to enterprise workflows and services.
\end{tcolorbox}

\section{Azure Implementation}

\begin{tcolorbox}[title={Azure AI Foundry}]
Use it as the central workspace for building, evaluating, deploying, and managing AI apps and agents on Azure.
It is suitable for enterprise GenAI projects that need models, agents, evaluations, safety checks, and deployment controls.
\end{tcolorbox}

\begin{tcolorbox}[title={Azure Machine Learning}]
Use it for training, experiment tracking, model registration, online endpoints, batch jobs, and MLOps.
It is a good fit for moving traditional ML and deep learning models from notebooks into governed production systems.
\end{tcolorbox}

\section{Databricks End-to-end AI Solutions}

\begin{tcolorbox}[title={Databricks Mosaic AI}]
Use it for building, evaluating, deploying, and governing GenAI and ML applications on the Databricks Lakehouse.
It fits roles where data engineering, model development, vector search, serving, and governance must live in one platform.
\end{tcolorbox}

\begin{tcolorbox}[title={Databricks Model Serving}]
Use it to deploy ML models, foundation models, and agents behind scalable serving endpoints.
It helps convert experiments into real-time or batch inference services that internal applications can call.
\end{tcolorbox}

\section{Data Pipelines and Feature Pipelines}

\begin{tcolorbox}[title={Apache Spark on Databricks}]
Use it for distributed data processing, feature generation, ETL, and large-scale transformations.
It is important when training data is large, comes from many systems, or must be processed regularly.
\end{tcolorbox}

\begin{tcolorbox}[title={Delta Lake}]
Use it as a reliable storage layer for versioned, auditable, and transactional data pipelines.
It helps keep training, inference, and analytics datasets consistent across teams and production jobs.
\end{tcolorbox}

\section{Training, Inference, Model Registry, and Monitoring}

\begin{tcolorbox}[title={MLflow}]
Use it to track experiments, package models, register model versions, and manage the model lifecycle.
For GenAI and agents, it also helps with tracing, evaluation, observability, prompt/version tracking, and production monitoring.
\end{tcolorbox}

\begin{tcolorbox}[title={Evidently AI}]
Use it to monitor data drift, prediction drift, model quality, and regression risks after deployment.
It is useful when production behavior may change because incoming data changes over time.
\end{tcolorbox}

\section{Performance, Scalability, and Maintainability}

\begin{tcolorbox}[title={Kubernetes / Azure Kubernetes Service}]
Use it to deploy, scale, and manage containerized AI services in production.
It is useful when multiple services, APIs, workers, and model endpoints must run reliably with autoscaling and rollout control.
\end{tcolorbox}

\begin{tcolorbox}[title={GitHub Actions}]
Use it for CI/CD: testing, linting, building Docker images, and deploying services after code review.
It supports maintainability because every change can be checked and released through a repeatable pipeline.
\end{tcolorbox}

\section{Modern AI Productivity Tools}

\begin{tcolorbox}[title={GitHub Copilot}]
Use it as an AI coding assistant for Python, tests, refactoring, documentation, and code exploration.
It can speed up implementation, but production code still needs review, security checks, and human ownership.
\end{tcolorbox}

\begin{tcolorbox}[title={Microsoft 365 Copilot}]
Use it to accelerate everyday knowledge work such as summarizing documents, drafting specifications, and preparing meeting notes.
In business-facing AI work, it can help translate stakeholder needs into clearer requirements and prototypes.
\end{tcolorbox}

\section{Collaboration with Business Units and Rapid Prototyping}

\begin{tcolorbox}[title={Streamlit}]
Use it to build quick internal demos, data apps, model dashboards, and proof-of-concept AI interfaces in Python.
It helps business users understand and test an AI use case before expensive production engineering begins.
\end{tcolorbox}

\begin{tcolorbox}[title={Power BI}]
Use it to communicate model outputs, KPI changes, and operational impact to non-technical stakeholders.
It is useful when AI results must be connected to reporting, decision-making, and business ownership.
\end{tcolorbox}

\section*{Suggested Interview Framing}
A strong answer for this role should not list tools randomly. A clear production architecture would be: data and features in Databricks, MLflow for experiment tracking and model registry, Azure AI Foundry/Azure OpenAI for GenAI components, Azure AI Search or Databricks vector search for RAG, FastAPI for internal APIs, Docker/Kubernetes for deployment, and MLflow or platform monitoring for quality, latency, drift, and cost.

\section*{Small Portfolio Project Idea}
Build an internal-document assistant: ingest PDFs or web pages, create embeddings, index them in Azure AI Search or Databricks vector search, answer questions with a RAG pipeline, expose the answer service through FastAPI, log traces and evaluation results with MLflow, and containerize the service with Docker. This single project demonstrates most requirements in the job advertisement.

\section*{Priorities if Time Is Limited}
\begin{enumerate}
    \item Azure AI Foundry / Azure OpenAI + Azure AI Search for enterprise GenAI and RAG.
    \item Databricks + Spark + Delta Lake + MLflow for data, ML lifecycle, and deployment.
    \item LangChain or LlamaIndex for RAG and LangGraph for agentic workflows.
    \item FastAPI + Docker + Kubernetes/Azure services for production integration.
    \item GitHub Copilot and Microsoft Copilot only as productivity tools, not as substitutes for engineering knowledge.
\end{enumerate}

\section*{Sources Consulted}
The recommendations are based on the job description and current official documentation or product pages for the tools below. Links are included for verification and further study.

\begin{itemize}
    \item Microsoft Foundry documentation: \url{https://learn.microsoft.com/en-us/azure/foundry/}
    \item Azure Machine Learning documentation: \url{https://learn.microsoft.com/en-us/azure/machine-learning/}
    \item Azure AI Search documentation: \url{https://learn.microsoft.com/en-us/azure/search/}
    \item Databricks Machine Learning documentation: \url{https://docs.databricks.com/en/machine-learning/}
    \item Databricks Model Serving documentation: \url{https://docs.databricks.com/en/machine-learning/model-serving/}
    \item LangChain documentation: \url{https://docs.langchain.com/}
    \item LlamaIndex documentation: \url{https://developers.llamaindex.ai/python/framework/}
    \item FastAPI documentation: \url{https://fastapi.tiangolo.com/}
    \item Kubernetes documentation: \url{https://kubernetes.io/docs/home/}
    \item MLflow documentation: \url{https://mlflow.org/docs/latest/}
    \item GitHub Copilot documentation: \url{https://docs.github.com/copilot}
    \item Streamlit documentation: \url{https://docs.streamlit.io/}
\end{itemize}

\bibliographystyle{plainnat}
\bibliography{/home/donkarlo/Dropbox/repo/nd_shared_research/bibliography/bibliography.bib}
\end{document}
